\chapter{Cronograma até a defesa}
\label{cap:cronograma}

\notaguia{Fonte: \texttt{Textos/backlog.md} (descontando o viés de fechamento)
+ conversa com o orientador. Os períodos da Quadro~\ref{quadro:cronograma} são
propostos e devem ser ajustados com a orientação. Item ainda a confirmar:
submissão de artigo, se o programa exigir produto técnico.}

O trabalho empírico está, em larga medida, concluído: as análises que
sustentam a tese estão executadas, auditadas e documentadas. O que resta é a
revisão de literatura --- a frente mais extensa e a menos comprimível ---,
o fechamento de arestas empíricas residuais e a redação final.

A revisão parte da agenda da Seção~\ref{sec:ref-agenda}: nove frentes, cada
uma com a pergunta que deve decidir. Não será um levantamento de cobertura, e
sim uma busca orientada por dúvidas específicas, em quatro etapas --- busca
estruturada nas bases (Web of Science, Scopus, SciELO e Google Scholar para a
produção regional em português), com registro das expressões e das datas;
triagem por título e resumo; leitura integral dos selecionados, que é o que
consome o tempo reservado no cronograma; e incorporação ao texto.

A prioridade é o levantamento dos trabalhos que testam deslocamento em escala
subnacional, do qual \citeonline{Arima2011} é o precedente conhecido, porque
dele depende a leitura de um resultado já obtido. As arestas empíricas
residuais estão inventariadas na Seção~\ref{sec:met-decisoes} e nenhuma
altera conclusão da tese. A redação final sucede, e não acompanha, o grosso
da revisão.

\begin{quadro}[htb]
\centering
\caption[Cronograma de atividades até a defesa]{Cronograma de atividades até a defesa (proposta a ajustar com a
orientação).}
\label{quadro:cronograma}
\small
\begin{tabular}{p{7.6cm}ccc}
\toprule
\textbf{Atividade} & \textbf{2026.2} & \textbf{2027.1} & \textbf{2027.2} \\
\midrule
Busca estruturada e triagem (agenda da Seção~\ref{sec:ref-agenda}) & $\bullet$ & & \\
\addlinespace
Leitura integral --- frentes prioritárias (teleacoplamento; composição da
expansão no Cerrado) & $\bullet$ & & \\
\addlinespace
Leitura integral --- demais frentes & & $\bullet$ & \\
\addlinespace
Reescrita do Capítulo~\ref{cap:referencial} a partir da revisão & & $\bullet$ & \\
\addlinespace
Fechamento das arestas empíricas residuais & & $\bullet$ & \\
\addlinespace
Revisão dos capítulos de resultados e discussão à luz da literatura & & $\bullet$ & $\bullet$ \\
\addlinespace
Redação final, revisão e formatação & & & $\bullet$ \\
\addlinespace
Depósito e defesa & & & $\bullet$ \\
\bottomrule
\end{tabular}
\fontefig{elaboração própria.}
\end{quadro}
